\documentclass[12pt,a4paper,psamsfonts]{amsart}
\usepackage{amssymb,mathtools}
\usepackage{calc}
\setlength{\topmargin}{0cm}
\setlength{\oddsidemargin}{0cm}
\setlength{\evensidemargin}{0cm}
\setlength{\marginparwidth}{0cm}
\setlength{\marginparsep}{0cm}
\setlength{\textheight}{\paperheight - 2in - 35pt}
\setlength{\textwidth}{\paperwidth - 2in}
\setlength{\headheight}{12.5pt}
\setlength{\headsep}{25pt}
\setlength{\footskip}{30pt}
\renewcommand{\baselinestretch}{1.3}
\pagestyle{headings}
\raggedbottom
\setlength{\emergencystretch}{2em}
\usepackage[colorlinks=true,urlcolor=blue,linkcolor=blue,citecolor=blue]{hyperref}
\newtheorem*{theorem}{Theorem}
\theoremstyle{definition}
\newtheorem*{setup}{Setup}
\theoremstyle{plain}
\newtheorem{lemma}{Lemma}[section]
\newtheorem{proposition}[lemma]{Proposition}
\newtheorem{corollary}[lemma]{Corollary}
\newcommand{\Q}{\mathbb Q}
\newcommand{\C}{\mathbb C}
\newcommand{\OO}{\mathcal O}
\newcommand{\mult}{\operatorname{mult}}
\newcommand{\nef}{\operatorname{n}}

\usepackage{tikz-cd}
\newcommand{\R}{\mathbb R}
\newcommand{\Z}{\mathbb Z}
\newcommand{\PP}{\mathbb P}
\DeclareMathOperator{\Pic}{Pic}
\DeclareMathOperator{\Supp}{Supp}
\DeclareMathOperator{\rk}{rk}
\DeclareMathOperator{\ord}{ord}
\DeclareMathOperator{\vol}{vol}
\DeclareMathOperator{\Sym}{Sym}
\DeclareMathOperator{\Bir}{Bir}
\DeclareMathOperator{\codim}{codim}
\def\UrlBreaks{\do\/\do-\do.}
\title{Nef Abundance: Framework -- Failure of Jet Separation in General Dimension}
\date{Temporary draft -- October 11, 2026}
\begin{document}
\maketitle

\section{Failure of jet separation}

Fix an integer $d\ge2$. We use the following source inputs,
denoted collectively by $(\mathsf I_d)$.
\begin{itemize}
\setlength{\itemsep}{2pt}
\setlength{\parsep}{0pt}
\item For every projective klt $\Q$-pair $(Y,B)$ of dimension
$j<d$ and every nef rational Cartier divisor $M$, if $K_Y+B$ is
pseudo-effective and, for some $a\in\Q_{>0}$,
\[
 (K_Y+B+aM)\cdot C>0
\]
for every integral curve through a sufficiently very general point,
then $K_Y+B+cM$ is big for every $c\in\Q_{>0}$.
\item Every projective log canonical $\Q$-pair of dimension at most
$d$ with nef adjoint has semiample adjoint.
\item Every projective $\Q$-factorial generalized klt $\Q$-pair of
dimension at most $d$, with fixed nef rational Cartier b-divisor
data and pseudo-effective adjoint, admits a terminating choice of
minimal-model program with a nef endpoint. The program extracts no
divisors, preserves $\Q$-factorial generalized klt models, and on a
common resolution the initial adjoint equals the final adjoint plus
an effective divisor exceptional over the final model.
\end{itemize}
The first input is the lower-dimensional assertion $(P_j)$ of
\cite[Section~3.1]{NS}. The last two are the companion inputs
used in \cite[Section~2]{NS}; here they are assumptions, not
independently established higher-dimensional theorems.

Let $X$ be a smooth projective variety over $\C$ of dimension $d$,
with $K_X\sim0$ and $q(X)=0$. Let $D$ be a nef Cartier divisor
such that $\nef(X,D)=d$ and $D^d=0$. Assume $(\mathsf I_d)$.

Assume in addition that $\kappa(D)=-\infty$.

\begin{setup}
\leavevmode
\begin{itemize}
\setlength{\itemsep}{2pt}
\setlength{\parsep}{0pt}
\item Fix an ample Cartier divisor $A$.
\item Fix $\epsilon\in\Q_{>0}$ sufficiently small. For integers
$n\ge1$, put
\[
 a_n=((D+A/n)^d)^{1/d},\qquad
 r_n=\max\{1,\lfloor\epsilon/a_n+1/2\rfloor\},\qquad
 P_n=r_n(D+A/n).
\]
Then $P_n$ is ample and rational Cartier, $r_n\to\infty$, and
$P_n^d\to\epsilon^d$. Indeed, $a_n\to0$ and
$|r_na_n-\epsilon|\le a_n/2$ for all sufficiently large $n$.
The symbol $n$ indexes the sequence; the dimension is $d$.
\item Let $p_i:X\times X\to X$ be the projections. Put
\[
 Z=\PP_{X\times X}
       \bigl(p_1^*\OO_X(D)\oplus p_2^*\OO_X(D)\bigr),
 \qquad \pi:Z\to X\times X.
\]
Let $S_1,S_2$ be the two summand sections and set
$U=Z\setminus(S_1\cup S_2)$. Thus $\dim Z=2d+1$.
\item Set $\xi=c_1(\OO_Z(1))$. For an integer $s>0$, put
\[
 R_{n,s}=(p_1\circ\pi)^*P_n+(p_2\circ\pi)^*P_n+s\xi.
\]
We use $\pi_*\OO_Z(a)=\operatorname{Sym}^a
(p_1^*\OO_X(D)\oplus p_2^*\OO_X(D))$ for $a\ge0$.
Whenever $|kR_{n,s}|$ is used, $k>0$ is an integer with
$kP_n$ Cartier.
\item Choose $\epsilon$ sufficiently small that
\[
 (14\epsilon)^d<\tfrac14,\qquad40\epsilon<\tfrac38.
\]
The integer $s$ is chosen after $\epsilon$ and before $n$ tends
to infinity. Any jet-separating multiple is chosen after $n$.
\end{itemize}
\end{setup}

\medskip
\noindent\textbf{Jet separation.}
Let $W$ be smooth, let $L$ be a Cartier divisor, and let $z\in W$.
For $j\ge0$, $|L|$ separates $j$-jets at $z$ if
\[
 \forall u\in\OO_{W,z}(L)\quad
 \exists\sigma\in H^0(W,\OO_W(L))\quad
 \sigma_z-u\in\mathfrak m_z^{j+1}\OO_{W,z}(L).
\]
\setcounter{lemma}{1}

\begin{lemma}\label{lem:general-nojets}
For every positive integer $s$ and every $n\gg0$,
$|kR_{n,s}|$ does not separate $(2d+2)k$-jets at any $z\in U$,
for any positive integer $k$ with $kP_n$ Cartier.
\end{lemma}

\medskip
\noindent\textbf{Draft status.}
This is a conditional, source-based proof framework following
\cite[Sections~4--6]{NS}. The source inputs $(\mathsf I_d)$ and
the inherited moving-centre and Frobenius arguments have not been
independently verified here. No unconditional higher-dimensional
nonvanishing or abundance theorem is asserted.
\section{Auxiliary estimates}
\label{sec:jet-tools}

The geometric argument will produce subvarieties with bounded degree and
bounded canonical intersection.  We first explain how these bounds
contradict a growing lower bound for the degrees of curves.  We then
record the moving-center results from \cite{LA} that supply the required
subvarieties.  All varieties in this section are over $\C$.

\subsection{The bounded-degree obstruction}

\begin{lemma}
\label{lemma:bounded-curves}
For every integer $f\geq1$ there is an integer $m_f>0$ with the
following property.  Let $Y$ be a smooth integral projective variety
of dimension $f$, let $D_0$ be a reduced simple normal crossing divisor
or the zero divisor, and let $L_0$ be a nef integral Cartier divisor
of full nef dimension on $Y$.  Suppose that $P_0$ is a nef rational
Cartier divisor and that $r,C_0,C_1$ are positive real numbers satisfying
\[
 0<P_0^f\leq C_0,\qquad P_0-rL_0\text{ is nef}.
\]
If $A_0=K_Y+D_0+L_0$ is big and
\begin{equation}
\label{j:adjoint-ratio}
 A_0\cdot P_0^{f-1}\leq C_1P_0^f,
\end{equation}
then
\begin{equation}
\label{j:bounded-curve-obstruction}
 r\leq(m_fC_1)^{f-1}C_0.
\end{equation}
The integer $m_f$ depends only on $f$, and not on the variety or
the divisors.  For dimensions $1\leq f\leq d$, one may replace all
$m_f$ by a single integer $m$ depending only on $d$.
\end{lemma}

\begin{proof}
The pair $(Y,D_0)$ is log canonical, its boundary coefficients belong
to $\{0,1\}$, and its nef polarizing divisor $L_0$ has Cartier index
one.  By \cite[Theorem~1.3]{BZ}, an integer $m_f$ depending only on
these data makes $|m_fA_0|$ birational.  Here $A_0$ is integral
Cartier, so no rounding or additional denominator is involved.

Resolve the base ideal by a projective birational morphism
$\mu:\widetilde Y\to Y$, with $\widetilde Y$ smooth, and write
\[
 \mu^*(m_fA_0)=H+F,
\]
where $H$ is the basepoint-free moving part and $F$ is effective.
The morphism defined by $H$ is birational onto its image; in
particular $H$ is nef and big.  In the rest of this proof, $P_0$
and $L_0$ denote their pullbacks to $\widetilde Y$.  Effectivity
of $F$ and nefness of $P_0$ give
\begin{equation}
\label{j:moving-adjoint-degree}
 H\cdot P_0^{f-1}\leq m_fC_1P_0^f.
\end{equation}

For $f\geq2$, put $s_i=P_0^{f-i}\cdot H^i$ for $0\leq i\leq f$.
Both divisors are nef and big, so the $s_i$ are positive.  The
Khovanskii--Teissier inequalities
\cite[Corollary~1.6.3(i) and Example~1.6.4]{LazarsfeldPAGI}
$s_i^2\geq s_{i-1}s_{i+1}$ show
that the ratios $s_i/s_{i-1}$ decrease.  Hence
\begin{equation}
\label{j:mixed-degree}
 P_0\cdot H^{f-1}
 \leq P_0^f\left(\frac{H\cdot P_0^{f-1}}{P_0^f}\right)^{f-1}
 \leq(m_fC_1)^{f-1}C_0.
\end{equation}
Choose a very general point $y\in\widetilde Y$ over the
isomorphism opens of both $\mu$ and the morphism defined by $H$.
We also choose $\mu(y)$ in the very general locus
that tests full nef dimension of $L_0$.  Pull back $f-1$ general
hyperplanes through the image of $y$ under the morphism defined
by $H$.  This subsystem has only the base point $y$, so the
general members meet properly and define an effective curve cycle
of class $H^{f-1}$; the component $C$ through $y$ is a curve
because their local equations are independent at $y$.
Nefness gives
\[
 P_0\cdot C\leq P_0\cdot H^{f-1}.
\]
The curve $C$ is not contracted by $\mu$, and its image passes
through the chosen very general point.  Full nef dimension and
integrality therefore give $L_0\cdot C\geq1$.  Since
$P_0-rL_0$ is nef, we obtain $P_0\cdot C\geq r$, which proves
\eqref{j:bounded-curve-obstruction}.

If $f=1$, take $C=Y$ directly.  Then
$P_0\cdot Y\geq rL_0\cdot Y\geq r$, and the required bound is
$r\leq C_0$.  Finally, \cite[Theorem~1.3]{BZ} applies to every
multiple of its prescribed integer.  A common multiple of the
finitely many integers $m_1,\ldots,m_d$ therefore gives the last
assertion.
\end{proof}

In particular, the hypotheses of Lemma~\ref{lemma:bounded-curves}
cannot hold along a sequence with $r\to\infty$ and with $C_0,C_1$
bounded, even if the varieties and divisors vary.  The following
observation explains how full nef dimension and pseudo-effectivity
of the canonical divisor pass to the moving subvarieties used below.

\begin{lemma}
\label{j:very-general-images}
Let $X$ be a normal integral projective variety that is not uniruled,
and let $L$ be a nef Cartier divisor of full nef dimension on $X$.
There is a complement $X^\circ$ of a countable union of proper
closed subsets of $X$ with the following property.  If $Y$ is a
smooth integral projective variety and $g:Y\to X$ is generically
finite onto a positive-dimensional image meeting $X^\circ$, then
$g^*L$ has full nef dimension, $Y$ is not uniruled, and $K_Y$ is
pseudo-effective.
\end{lemma}

\begin{proof}
Choose $X^\circ$ so that every curve through any of its points has
positive $L$-degree and no rational curve meets $X^\circ$.  The first
condition is the very general characterization of full nef
dimension.  For the second, rational curves are parametrized by
countably many algebraic families; the closure of the evaluation
image of each such family is proper because $X$ is not uniruled.

Since $g(Y)$ meets $X^\circ$, it is contained in none of the closed
sets excluded in its definition.  A very general point of $Y$ thus
maps into $X^\circ$ and lies in the quasi-finite locus of $g$.
Every curve through such a point has a curve as its image, so its
$g^*L$-degree is positive.  This proves full nef dimension.  A
covering family of rational curves on $Y$ would similarly give a
nonconstant rational curve meeting $X^\circ$, a contradiction.
The pseudo-effectivity of $K_Y$ now follows from \cite{BDPP}.
\end{proof}

\subsection{Subadjunction and high-multiplicity components}

We use the following five results from the section
\emph{Tools for moving base components} of \cite{LA}.  Their
hypotheses distinguish generic information along a center from
global information about the ambient variety.  Intersections on
an integral subvariety are computed after pullback to its
normalization and a smooth projective resolution.  A base ideal
is \emph{isolated at the generic point} of a subvariety when its
localization there is primary for the maximal ideal.

\begin{lemma}
\label{j:subadjunction}
Let $Z$ be a projective $\Q$-factorial klt variety, let
$\Theta\geq0$ be a rational divisor, and let $V\subset Z$ be
integral of dimension $f>0$.  Suppose that $(Z,\Theta)$ is log
canonical at the generic point of $V$ and has a divisor of log
discrepancy zero with center $V$.  If $P$ is a nef rational Cartier
divisor on $V$ and $\rho:V^*\to V$ is a smooth projective resolution
factoring through its normalization, then
\[
 K_{V^*}\cdot(\rho^*P)^{f-1}
 \leq (K_Z+\Theta)|_V\cdot P^{f-1}.
\]
\end{lemma}

Only generic log canonicity is required.  The proof in \cite{LA}
uses the dlt modification of \cite[Theorem~2.10]{FH}, whose
truncated boundary leaves an effective surplus supported over
the non-lc locus.  On a minimal stratum over $V$, adjunction
gives a pair that is klt over the generic point of its
connected-fiber base.  This suffices for the klt-trivial
canonical bundle formula and its b-nef moduli divisor
\cite[Definition~3.1 and Theorem~3.5]{FG}.  The effective
discriminant, the moduli divisor, and finite ramification have
nonnegative intersections with the pulled-back nef class.
The comparison uses canonical Weil cycles, so it requires no
$\Q$-Gorenstein assumption on the normalization of $V$.

\begin{lemma}
\label{j:base-component}
Let $Z$ be a projective $\Q$-factorial klt variety of dimension
$e$, let $R$ be a nef rational Cartier divisor, and let $k>0$
be an integer such that $kR$ is Cartier.  Let
$0\ne\mathcal H\subset H^0(Z,\OO_Z(kR))$ be a linear subspace
with proper base locus and base ideal $\mathfrak b$.  Suppose that
$V$ is an integral base-locus component of dimension $f<e$,
isolated at its generic point $\eta_V$, and that
$\eta_V\in Z_{\mathrm{sm}}$.  Let
$\mathfrak m$ be the maximal ideal of $\OO_{Z,\eta_V}$.
If $\mathfrak b\OO_{Z,\eta_V}\subset\mathfrak m^h$ for an
integer $h>0$, there are a rational number $0<c\leq(e-f)k/h$ and an
effective rational divisor $\Theta\sim_{\Q}cR$ such that
$(Z,\Theta)$ is log canonical at $\eta_V$ and has a log canonical
place centered on $V$.  The degree and canonical bounds are
\begin{equation}
\label{j:base-bounds}
 \begin{split}
 R^f\cdot V&\leq(k/h)^{e-f}R^e,\\
 K_{V^*}\cdot R^{f-1}
   &\leq(K_Z+cR)|_V\cdot R^{f-1}\qquad(f>0),
 \end{split}
\end{equation}
where $V^*$ is a smooth projective resolution.  The degree bound
also holds for $f=0$.
\end{lemma}

The degree estimate is the local multiplicity bound $h^{e-f}$
combined with successive general cuts by the given linear series.
The boundary is obtained by averaging general members with total
coefficient equal to the generic log canonical threshold of
$\mathfrak b$.  Thus it is an actual divisor to which
Lemma~\ref{j:subadjunction} applies.

\subsection{Full jet kernels and restriction to a fiber}

\begin{lemma}
\label{j:moving-kernel}
Let $Z$ be an integral projective variety of dimension $e$, let
$R$ be a nef, big, semiample rational Cartier divisor, and fix
an integer $k>0$ such that $kR$ is Cartier.  Choose positive real
numbers $\tau_0<\cdots<\tau_e$ with consecutive gap $\delta>0$.
For $x\in Z_{\mathrm{sm}}$, put
\[
 \mathcal H_i(x)=H^0\bigl(Z,\OO_Z(kR)
                  \otimes\mathfrak m_x^{\lceil k\tau_i\rceil}\bigr).
\]
Suppose that for very general $x$ all these spaces are nonzero and
the base locus of $\mathcal H_0(x)$ has a positive-dimensional
component through $x$.  If $k\delta\geq4$, then, after an
algebraic parameter extension and restriction to nonempty opens,
there are an integral parameter variety $T$ with a dominant
marked-point map $x:T\to Z_{\mathrm{sm}}$, a fixed index
$i\in\{0,\ldots,e-1\}$, and a family of integral subvarieties
$V_t$ containing $x(t)$ and common to the base loci of
$\mathcal H_i(x(t))$ and $\mathcal H_{i+1}(x(t))$.  In particular,
the incidence
\[
 \Gamma=\{(t,z):z\in V_t\}\subset T\times Z
\]
dominates $Z$.

For general $t$, the dimension $f$ of $V_t$ lies in $[1,e-1]$,
$V_t$ is an isolated component of the higher base locus at its
generic point, and the entire higher-kernel base ideal is
contained there in the ordinary power $\mathcal I_{V_t}^h$, where
\[
 h=\lceil k\delta/2\rceil,\qquad k/h\leq2/\delta.
\]
The incidence base ideal is likewise contained in
$\mathcal I_\Gamma^h$ on a dense open of $\Gamma$.
These containments are taken in the smooth ambient open.
If $Z$ is also $\Q$-factorial and klt, then
\begin{equation}
\label{j:moving-bounds}
 \begin{split}
 R^f\cdot V_t&\leq(2/\delta)^{e-f}R^e,\\
 K_{V_t^*}\cdot R^{f-1}
   &\leq(K_Z+cR)|_{V_t}\cdot R^{f-1},
       \qquad 0<c\leq2(e-f)/\delta.
 \end{split}
\end{equation}
The top intersection on $V_t$ is positive when $V_t$ meets the
open on which the big semiample contraction is an isomorphism.
\end{lemma}

The use of the full kernels matters: differentiation in the
parameter lowers vanishing at the moving mark and still gives a
global section in the lower kernel.  Incidence dominance makes
these parameter directions span the normal directions.  The
vanishing of all derivatives of order less than $h$ then gives
the ordinary-power containment, which can be restricted to
general fibers as follows.

\begin{lemma}
\label{j:fiber-restriction}
In the setting of Lemma~\ref{j:moving-kernel}, let $g:Z\to B$ be
a morphism.  There is a dense open of the incidence with the
following property.  Choose $(t,z)$ in this open, put $b=g(z)$,
and suppose that $V_t$ is smooth over a smooth open of its actual
image at $z$.  Let $F$ be the reduced component through $z$ of
the scheme fiber of $V_t\to g(V_t)$.  Near $z$ in $Z_b$, the
restricted higher-kernel base ideal has support exactly $F$ and
is contained in the ordinary power $\mathcal I_F^h$.

If $F$ is proper in an integral ambient fiber component containing
it, the restricted series is nonzero there and $F$ is an isolated
base component at its generic point.  Applying
Lemma~\ref{j:base-component} on that ambient component additionally
requires its projectivity, $\Q$-factoriality, klt singularities,
and smoothness at the generic point of $F$.
\end{lemma}

Here the fiber is chosen through a general incidence point.  The
statement does not assert the same conclusion for an arbitrary
special fiber.  Smoothness over the actual image identifies the
restricted ideal of $V_t$ with the reduced ideal of $F$ locally;
this is what preserves ordinary powers.

The $\Q$-factoriality needed for the base-component estimates can
also be supplied by a small projective $\Q$-factorialization
$\sigma:Z'\to Z$ of a klt ambient variety.  This passage is made
only for components that meet the isomorphism open of $\sigma$.
Their strict transforms retain the generic base-ideal and
multiplicity conditions for the pulled-back series, and the
testing class $\sigma^*R$ is nef.  Since
$K_{Z'}=\sigma^*K_Z$, projection formula gives the same degree
and canonical-intersection bounds on a common resolution.
The pullback testing class need not be ample; nefness is the
property used in these estimates.  Passing to higher resolutions
does not change the canonical intersections, because their
additional exceptional divisors have zero intersection with
the pulled-back testing class to the required power.

\subsection{Finite covers with a fixed branch complement}

A bound on the degree of a branched cover does not make the
possible covers finite.  The final companion input supplies the
needed fixed branch complement from a family whose branch
divisors do not sweep the base.

\begin{lemma}
\label{j:finite-covers}
Let $Y$ be a normal integral projective variety, let
$\mathcal V\to T$ be a smooth projective family with integral
fibers over an integral parameter variety, and let
$g:\mathcal V\to Y$ be a morphism.  Suppose that the general
fiber maps $g_t:V_t\to Y$ are dominant and generically finite of
degree at most a fixed integer $e$.

After a finite parameter extension and shrinking, spread the
geometric generic components of the zero divisor of the relative
Jacobian over $Y_{\mathrm{sm}}$ to divisors
$\mathcal R_1,\ldots,\mathcal R_N$ in $\mathcal V$.
Retain this notation after the base change.  Let $I$ consist of
the indices $j$ for which
$\overline{g_t((\mathcal R_j)_t)}$ is a divisor in $Y$ for general
$t$.  Assume that
\[
 \overline{g(\mathcal R_j)}\ne Y\qquad(j\in I).
\]
After shrinking again, there is a smooth nonempty open
$Y^\circ\subset Y$ over which every general finite normal Stein
cover of $g_t$ is finite \'etale.  Only finitely many such normal
covers of $Y$ occur, up to isomorphism over $Y$.
The conclusion holds simultaneously for finitely many
projections; the open and the finite lists may depend on all
fixed parameters of the family.
\end{lemma}

To see the role of the hypothesis, remove $Y_{\mathrm{sing}}$
and the finitely many proper image closures
$\overline{g(\mathcal R_j)}$ for $j\in I$.  Every branch prime
of a general Stein cover is accounted for by one of these
relative Jacobian components.  Purity therefore makes the
restricted covers \'etale.  The fundamental group of the
resulting smooth complex variety is finitely generated, so it
has only finitely many covers of degree at most $e$.  Finally,
a finite normal cover of $Y$ is determined by its function
field, hence by its restriction to this fixed open.

\section{Geometric estimates}\label{sec:geometry}
In the technical arguments below, put $L=D$, $t=1/n$, $r=r_n$,
$P=P_n$, and $q=s$. The letter $L$ keeps the source notation for
the fixed nef Cartier divisor; $q$ is the bundle parameter.
\begin{equation}\label{eq:scaling}
 P=r(L+tA),\qquad r\to\infty,\qquad P^d\to\epsilon^d.
\end{equation}
For each fixed $q$, discard finitely many terms so that
\begin{equation}\label{j:volume-window}
 \epsilon^d/2\le P^d\le2\epsilon^d,\qquad r>q+1.
\end{equation}
Every curve through a very general point of $X$ has positive integral
$L$-degree, hence $P$-degree at least $r$. Thus a covering family of
curves of bounded $P$-degree will give a contradiction.

We specify how the same obstruction applies to subvarieties and
covers. A resolution $W\to X$ has an effective canonical divisor:
pull back a nowhere-zero canonical form on the smooth locus of $X$
and use terminality. Hence $W$, and therefore $X$, is not uniruled
\cite{BDPP}. By Lemma~\ref{j:very-general-images}, if a smooth
projective variety $Y$ maps generically finitely to a
positive-dimensional subvariety containing a sufficiently very general
point of $X$, then $Y$ is not uniruled and the pulled-back Cartier
divisor $L_Y$ has full nef dimension. In particular, $K_Y$ is
pseudo-effective. If $f=\dim Y<d$, the induction hypothesis $(P_f)$
applied to $(Y,0)$ gives
\begin{equation}\label{j:lower-big}
                         K_Y+L_Y\ \text{big}.
\end{equation}
All occurrences of ``very general'' below include the fixed
conditions of Lemma~\ref{j:very-general-images}.

Consequently, for $f<d$ it suffices to find on such a $Y$ a nef
testing divisor $P_Y$ with
\[
 0<P_Y^f\le C_0,\qquad P_Y-rL_Y\ \text{nef},\qquad
 K_YP_Y^{f-1}\le C P_Y^f,
\]
where $C_0,C$ are independent of $t$. Indeed
$L_YP_Y^{f-1}\le P_Y^f/r$, so Lemma~\ref{lemma:bounded-curves}
applies. We will also use that lemma in dimension $d$, after proving
the required bigness separately. For the rest of this section, a
bound is called uniform if it is independent of $t$; it may depend
on parameters fixed before $t$ tends to zero. Since the dimensions
lie in a finite range, Lemma~\ref{lemma:bounded-curves} excludes all
such configurations once $r$ is sufficiently large.

\subsection{The scalar jet estimate}

Put $\rho=(P^d)^{1/d}$. The first estimate bounds the vanishing order
of every section at a very general point; it is stronger than a
bound on a single fixed linear system.

\begin{lemma}\label{lemma:scalar-bound}
In the setting of \eqref{eq:scaling}, for all sufficiently small $t$
and at a very general point $x\in X$, every nonzero section of $kP$,
where $k$ is a positive integer and $kP$ is Cartier, satisfies
\begin{equation}\label{j:scalar-order}
                         \ord_x(s)\le4k\rho.
\end{equation}
\end{lemma}

\begin{proof}
Suppose that violations occur for arbitrarily small $t$. Fix such a
$t$. At a very general point a violation for some $k$ implies that
the corresponding complete jet kernel is nonzero at the general
point. Indeed the kernel ranks are algebraic rank conditions on the
smooth locus, and there are only countably many choices of $k$ and
of the required integral order. Raising the violating section to
powers allows $k$ to be arbitrarily large and divisible.

Choose $d+1$ levels
\[
 \tau_i=\frac32\rho+i\delta\quad(0\le i\le d),
 \qquad \delta=\frac{2\rho}{d}.
\]
They all lie strictly between $\rho$ and $4\rho$. For a sufficiently
large divisible $k$, their full jet kernels are nonzero and
$k\delta\ge4$. The lowest kernel has a positive-dimensional base
component through its marked point. Otherwise that point would be
an isolated base component of multiplicity at least
$h_0=\lceil k\tau_0\rceil$. The zero-dimensional case of
Lemma~\ref{j:base-component} would give
\[
                1\le (k/h_0)^nP^d<1,
\]
a contradiction.

Lemma~\ref{j:moving-kernel} now supplies a moving family of proper
base components $V$, with dominant incidence and dimension
$1\le f<d$. On a resolution $V^*$, write $P_V,L_V$ for the
restrictions followed by pullback. The lemma gives
\[
 0<P_V^f\le(2/\delta)^{d-f}P^d,
 \qquad K_{V^*}P_V^{f-1}\le cP_V^f,
 \qquad 0<c\le\frac{2(d-f)}{\delta}.
\]
Here the canonical term of the ambient variety vanishes because
$K_X\sim0$. The volume window \eqref{j:volume-window} bounds
$\delta$ above and away from zero, so these are uniform estimates.
The incidence point may be chosen over a very general point of $X$.
Thus \eqref{j:lower-big} gives bigness of $K_{V^*}+L_V$, and
$P_V-rL_V$ is nef. Lemma~\ref{lemma:bounded-curves} contradicts
$r\to\infty$.
\end{proof}


\section{The fixed witness}
\begin{lemma}\label{lem:general-witness}
For every fixed $s>0$ and every $n\gg0$, there is a very general
point $x\in X$ such that, on its blowup
$\beta:\widetilde X\to X$ with exceptional divisor $E$, there
are a birational morphism $f:V\to\widetilde X$ from a smooth
projective variety and ample Cartier divisors $H_1,\ldots,H_{d-1}$
on $V$ satisfying
\[
 f^*\beta^*(P_n+sD)\cdot H_1\cdots H_{d-1}
 <40\epsilon\, f^*E\cdot H_1\cdots H_{d-1}.
\]
\end{lemma}
\begin{proof}
Put $\rho_n=(P_n^d)^{1/d}$ and choose $n$ sufficiently large that
Lemma~\ref{lemma:scalar-bound} applies, $r_n>s$, and $\rho_n<2\epsilon$.
Every nonzero section of $2kP_n$ then has order at most
$8k\rho_n<16k\epsilon$ at a very general point $x$.
The divisor
\[
 \beta^*(2P_n)-40\epsilon E
\]
is not pseudo-effective. Otherwise its convex combination with
the big divisor $\beta^*(2P_n)$ would make
$\beta^*(2P_n)-cE$ big for some rational
$16\epsilon<c<40\epsilon$. A sufficiently divisible multiple
would have a section of order at least $kc$, a contradiction.

By pseudo-effective/strongly-movable duality \cite{BDPP}, choose
one actual generating complete intersection on a smooth birational
model whose pairing with this divisor is negative. Approximate
the ample real classes by rational ones and clear denominators
to obtain $f,H_1,\ldots,H_{d-1}$ with
\[
 f^*\beta^*(2P_n)\cdot H_1\cdots H_{d-1}
 <40\epsilon\,f^*E\cdot H_1\cdots H_{d-1}.
\]
Since $P_n-sD$ is nef, the required inequality follows.
The denominator is positive because $f^*E$ is nonzero effective
and all the $H_i$ are ample. Fix this model and these divisors
before reduction.
\end{proof}

\begin{lemma}\label{lem:general-flag}
For every ample rational Cartier divisor $P$ on $X$, there are
positive integers $l_1,\ldots,l_{d-1}$ and a smooth integral
complete-intersection flag of successive classes $l_iP$,
ending in a curve $C$, such that
\[
 \mu_{\min}(\Omega_X^1|_C)\ge0.
\]
\end{lemma}
\begin{proof}
Since $K_X\sim0$, $X$ is not uniruled. Generic semipositivity of
$\Omega_X^1$ \cite[Theorem~2.1]{CP} and restriction of its
Harder--Narasimhan filtration to sufficiently general complete
intersections \cite[Theorem~6.1]{MR} give the asserted slope
bound. Choose the $l_i$ sufficiently large and divisible that
$l_iP$ is very ample; Bertini gives the smooth integral flag.
\end{proof}

\section{The Frobenius obstruction}\label{sec:frobenius}
The finite-data statement below is \cite[Theorem~6.1]{NS},
specializing the companion Frobenius comparison. Its divisor $G$
is a rational polarization and is distinct from the fixed nef
Cartier divisor $D$ in the Setup. All models, sections, the flag,
and the complete-intersection witness are fixed before reduction.
\subsection{The finite-data statement}

For a vector bundle on a smooth projective curve, its slope is
degree divided by rank. The least and greatest Harder--Narasimhan
slopes are denoted by $\mu_{\min}$ and $\mu_{\max}$.
For a line bundle on a smooth variety, to generate jets through
order $a$ at a point means surjectivity to its stalk modulo
the $(a+1)$-st power of the maximal ideal.

\begin{proposition}
\label{thm:finite-data}
Let $W$ be a smooth connected projective complex variety of
dimension $d\ge2$. Let $G,H$ be rational Cartier divisors, with
$H$ ample, let $S$ be an integral Cartier divisor, and fix an
integer $q>0$ and real numbers $r>1$, $\epsilon>0$. Suppose
\begin{equation}\label{fr:intersection-bounds}
 \begin{gathered}
 H^d\le(2\epsilon)^d,\qquad
 K_WH^{d-1}\le2H^d/r,\qquad
 (2G+qS)H^{d-1}\le4H^d,\\
 (14\epsilon)^d<\tfrac14,\qquad 80\epsilon<\tfrac34.
 \end{gathered}
\end{equation}
Assume that the following fixed data exist.
\begin{enumerate}
\renewcommand{\labelenumi}{\textup{(\roman{enumi})}}
\item A smooth integral complete-intersection flag from $W$ to
a curve $C$, whose successive divisor classes are $h_i=l_iH$
for fixed positive integers $l_i$ for which $l_iH$ is Cartier, with
$\mu_{\min}(\Omega_W^1|_C)\ge0$.
\item On the projective bundle
\[
 Z=\PP\bigl(\mathcal O_W(S)_1\oplus
                         \mathcal O_W(S)_2\bigr)\longrightarrow W\times W,
 \qquad \xi=c_1(\mathcal O_Z(1)),\qquad M=G_1+G_2+q\xi,
\]
use the convention that the direct image of $\mathcal O_Z(a\xi)$
is the $a$-th symmetric power of the displayed sum for $a\ge0$;
subscripts indicate pullback from the corresponding factor of
$W\times W$. The two axes are the sections defined by the two
summands. A smooth point $z_*$ in their complement and finitely many
sections of $k_0M$, for some integer $k_0>0$ with $k_0G$ Cartier,
generate jets through order $(2d+2)k_0$ at $z_*$.
\item A point $x\in W$, its blowup
$\pi_x:\widehat W\to W$ with exceptional divisor $J$, and a
curve class
\[
 \gamma=f_*(A_1\cdots A_{d-1}),
\]
where $f:V\to\widehat W$ is one fixed birational morphism,
$V$ is smooth integral projective, and the $A_i$ are ample
integral Cartier divisors, such that $J\cdot\gamma>0$ and
\begin{equation}\label{fr:movable-input}
 \pi_x^*(G+K_W/2+\lambda S)\cdot\gamma
 \le40\epsilon J\cdot\gamma\qquad(0\le\lambda\le q).
\end{equation}
It suffices to check the two endpoints in this affine inequality.
\end{enumerate}
These data cannot coexist.
\end{proposition}

The point $x$ is independent of the two base coordinates of $z_*$.
No nefness is required of $G,S$, or $K_W$, and the theorem does not
assume that $K_W$ is pseudo-effective. Its hypotheses specify the
flag and the movable curve class themselves. All these data will
be fixed before the residual characteristic tends to infinity.

\subsection{Reduction and Frobenius jets}
\label{fr:reduction}

We begin the proof of Proposition~\ref{thm:finite-data} by spreading
its hypotheses and converting the fixed ordinary jets
into jets of order proportional to the residual characteristic.
Put
$\Lambda=\prod_{i=1}^{d-1}l_i$. The flag satisfies
\begin{equation}\label{fr:flag-degrees}
 [C]=\Lambda H^{d-1},\qquad h_i\cdot C=l_i\Lambda H^d.
\end{equation}
Choose one sufficiently ample integral Cartier divisor $T_0$ so
that every $T_0+aS$, for integers $0\le a\le(k_0-1)q$, has a section
nonvanishing at each of the two base coordinates of $z_*$.
Fix these finitely many sections. All models, morphisms, divisors,
points, the flag, the complete-intersection witness for $\gamma$,
and the jet and filler sections spread over an integral finitely
generated $\Z$-algebra of characteristic zero. Shrink its
spectrum so that the fibers and flag are smooth projective and
geometrically integral, the birational morphisms remain birational,
the fixed ample bundles remain ample, and the finite jet
surjection and nonvanishings persist. Birationality is preserved
by spreading an isomorphism between dense opens. All displayed
intersection numbers are constant.

We also preserve $\mu_{\min}(\Omega_W^1|_C)\ge0$.
One way is to spread its characteristic-zero Harder--Narasimhan
filtration, make all its factors locally free, and use openness
of their semistability on curves; their degrees and ranks are
fixed. Equivalently, a single relative ample twist globally
generates the curve bundle. Every quotient of rank $a$ then has
degree bounded below by $-ad$ for a fixed integer $e$.
A negative-degree quotient therefore has one of finitely many ranks
and degrees. The proper relative Quot schemes for those Hilbert
polynomials have images missing the generic point: after removal
of torsion, a negative-degree coherent quotient there would
give a negative-degree vector-bundle quotient. Removing those
finitely many images proves the same openness assertion directly.

This base has closed points in arbitrarily large characteristics.
Take algebraic closures of their residue fields and exclude
denominator primes. Retain the notation for the reductions.
The curve $\gamma$ continues to pair nonnegatively with every
effective divisor: pull that divisor back under the fixed
birational morphism and intersect with its fixed ample divisors.
This tests effective divisors on the reduction itself. No
specialization of a characteristic-zero effective cone is used.

\begin{lemma}\label{fr:positive-characteristic-jets}
For a sufficiently large residual characteristic $p$, set
\[
 N_p=\lfloor p/k_0\rfloor,\qquad B_p=k_0N_pD+T_0.
\]
The sections of $pq\xi+B_{p,1}+B_{p,2}$ generate jets through order
$(2d+2)k_0N_p$ at $z_*$.  In particular, they surject onto the quotient
of its local ring by the $p$-th powers of all regular parameters.
\end{lemma}

\begin{proof}
Multiply $N_p$ copies of the fixed jet system.  Every monomial of
total degree at most $(2d+2)k_0N_p$ is a product of $N_p$ monomials
of degree at most $(2d+2)k_0$.  Taking sections with those leading
monomials gives a triangular system, ordered by total degree.
Starting at degree zero and removing the error in each successive
degree proves the required jet surjection. This argument takes
place in the local ring and never divides by factorials.

Put $e_p=(p-k_0N_p)q$, so $0\le e_p\le(k_0-1)q$.
Among the sections fixed before reduction, choose
\[
 u_0\in H^0(W,\OO_W(T_0)),\qquad
 v_{e_p}\in H^0(W,\OO_W(T_0+e_pS))
\]
nonzero at the first and second base coordinates of $z_*$,
respectively. Their tensor product belongs to the summand of
first-slot weight zero in the projective-bundle formula; its
fiber monomial is the $e_p$-th power of the second coordinate.
It defines a section of
$e_p\xi+T_{0,1}+T_{0,2}$ nonzero at $z_*$, since $z_*$ lies
off both axes. Multiplying the product jet system by this section
changes its divisor from $N_pk_0M$ to
$pq\xi+B_{p,1}+B_{p,2}$. Multiplication is invertible on the
local jet algebra. Only the finitely many previously fixed
values of $e_p$ occur.

Since $\dim Z=2d+1$, the
quotient by the $p$-th powers of regular parameters has top total
degree $(2d+1)(p-1)$.  For large $p$ this is at most
$(2d+2)k_0N_p$. The underlying ordinary/Frobenius ideal comparison is
also recorded in \cite[proof of Proposition~2.12, Equation~(2.8)]{MustataSchwede}.
\end{proof}

\subsection{Full rank from the two-slot jets}
\label{fr:full-rank}

Let $F:W\to W'$ be relative Frobenius to the base-field twist.
It is finite flat because $W$ is smooth.  Write $S'$ for the twist
of $S$, so $F^*S'=pS$, and define
\[
 \mathcal E=F_*\OO_W(B_p),\qquad s=\rk\mathcal E=p^d,
 \qquad R=s^2.
\]
For integers $0\le a\le q$, put
$U_a=H^0(W,\OO_W(B_p+paS))$.
Projection formula gives an evaluation map on $W'\times W'$:
\begin{equation}\label{fr:two-slot-evaluation}
 \bigoplus_{\substack{a+b=q\\a,b\ge0}}
 U_a\otimes U_b\otimes
 \OO(-aS'_1-bS'_2)
 \longrightarrow \mathcal E\boxtimes\mathcal E.
\end{equation}

\begin{lemma}\label{fr:generic-rank}
The map in \eqref{fr:two-slot-evaluation} has generic rank $R$.
\end{lemma}

\begin{proof}
Trivialize the two summands defining $Z$ near the selected base
points.  Let $z$ be the torus coordinate, whose value at $z_*$ is
$z_0\ne0$.  By the projective-bundle formula, the sections in
Lemma~\ref{fr:positive-characteristic-jets} are weight polynomials in $z$
with weight-$j$ coefficients in
\[
 H^0(W,B_p+jS)\otimes
 H^0(W,B_p+(pq-j)S),\qquad 0\le j\le pq.
\]
Let $A_*$ be the tensor product of the local algebras of the two
base Frobenius fibers.  The local quotient in
Lemma~\ref{fr:positive-characteristic-jets} is
\[
 A_*[z]/(z^p-z_0^p),
\]
free over $A_*$ with basis $1,z,\ldots,z^{p-1}$.  Project to the
coefficient of $1$ in this basis.  Exactly the weights $j=pa$
survive, with multipliers $z_0^{pa}$.  Their coefficients are the
values of the corresponding summands of
\eqref{fr:two-slot-evaluation}, in the chosen frames and the induced
twisted frames.  As the full jet map is surjective, these coefficients
span $A_*$, of dimension $p^{2d}=R$.  Thus the evaluation has full
rank at this pair of base points, and hence generically.  The
calculation uses a basis over the possibly nonreduced algebra $A_*$,
so reducedness of the Frobenius fibers is unnecessary.
\end{proof}

The jet calculation at $z_*$ proves generic full rank; the diagonal
estimate below will be tested at $(x',x')$, obtained from the
separately fixed point $x$.

\subsection{Sections on the Frobenius diagonal}
\label{fr:filtration}

The Frobenius fiber product has the cartesian square
\begin{equation}\label{fr:frobenius-square}
\begin{tikzcd}
 \Delta_F=W\times_{W'}W \arrow[r,"\mathrm{pr}_2"] \arrow[d,"\mathrm{pr}_1"'] & W \arrow[d,"F"]\\
 W \arrow[r,"F"'] & W'.
\end{tikzcd}
\end{equation}
Its reduced subscheme is the diagonal $W$.  Put
\[
 \mathcal L_p=
 \OO_{\Delta_F}(B_{p,1}+B_{p,2}+pqS_1).
\]
Write $h=F\circ\mathrm{pr}_1=F\circ\mathrm{pr}_2:
\Delta_F\to W'$. On this fiber product,
\[
 \OO(pS_1)\simeq h^*\OO_{W'}(S')\simeq\OO(pS_2).
\]
Consequently every pair of weights $a+b=q$ gives the same line bundle:
\[
 \OO_{\Delta_F}(B_{p,1}+B_{p,2}+paS_1+pbS_2)
 \simeq\mathcal L_p.
\]
On the diagonal of $W'\times W'$, every source twist of
\eqref{fr:two-slot-evaluation} becomes $-qS'$.
Finite base change and projection formula identify the twisted target as
\[
 (\mathcal E\otimes\mathcal E)(qS')\simeq h_*\mathcal L_p.
\]
The product sections defining the columns restrict to global sections
of this one bundle $\mathcal L_p$. Their values at a diagonal point
therefore lie in the image of $H^0(\Delta_F,\mathcal L_p)$ in the
corresponding fiber of $h_*\mathcal L_p$. This proves that the rank
at every diagonal point is at most
\begin{equation}\label{fr:diagonal-rank-upper}
 h^0(\Delta_F,\mathcal L_p).
\end{equation}

For a rank-$d$ bundle $V$ in characteristic $p$, denote by $A_j(V)$
the degree-$j$ part of its symmetric algebra modulo the $p$-th powers
of local generators.  This construction is independent of frame:
the $p$-th power of a linear combination is the sum of the $p$-th
powers in characteristic $p$.  Each $A_j(V)$ is locally free and a
quotient of $V^{\otimes j}$.  Set
\[
 u=d(p-1),\qquad e_j=\rk A_j(V),\qquad
 \sum_{j=0}^u e_j=p^d.
\]
Filtering $\mathcal L_p$ by powers of the ideal of the reduced
diagonal in $\Delta_F$ gives the graded bundles
\begin{equation}\label{fr:diagonal-graded}
 \mathcal G_j=\OO_W(2B_p+pqS)\otimes A_j(\Omega_W^1),
 \qquad 0\le j\le u.
\end{equation}
Indeed, in smooth local coordinates the algebra is generated by
parameter differences with their $p$-th powers zero, and its
degree-one conormal is $\Omega_W^1$.  This coordinate calculation
can be made étale locally or in completions and identifies the
global associated graded.  This is the diagonal-ideal form of the canonical Frobenius filtration
\cite[Corollary~3.5]{KS}. The same calculation, using
only the bundle from the second factor, filters $F^*\mathcal E$ with pieces
\cite[Theorem~3.7]{Sun}
\begin{equation}\label{fr:Frob-pullback-graded}
 \OO_W(B_p)\otimes A_j(\Omega_W^1).
\end{equation}

Complementary multiplication is a perfect pairing, giving
\begin{equation}\label{fr:socle-pairing}
 A_{u-j}(V)\simeq A_j(V)^\vee\otimes(\det V)^{p-1}.
\end{equation}
On monomials, each exponent vector pairs with its complement to
$(p-1,\ldots,p-1)$.  The top line transforms by the character
$(\det V)^{p-1}$: this is immediate on diagonal matrices and hence
 identifies the character of $\mathrm{GL}_d$ on that line.

\subsection{A uniform bound for the diagonal rank}
\label{fr:rank-bound}

We bound the sections of every graded bundle in
\eqref{fr:diagonal-graded} by its rank times the same scalar
polynomial in $p$. This uniformity lets us sum the ranks of the
graded pieces without an additional factor for their number.
We begin with a slope bound on the fixed curve; here slope means
ordinary degree divided by rank.

\begin{lemma}\label{fr:truncated-slope}
There is a constant $c\ge0$, independent of the sufficiently large
residual characteristic $p$ and of $j$, such that for
$V=\Omega_W^1|_C$ one has
\[
 \mu_{\max}(A_j(V))\le(p-1)K_W\cdot C+nc.
\]
\end{lemma}

\begin{proof}
Write $F_C$ for absolute Frobenius of $C$.  Langer's instability
estimate \cite[Corollary~2.5]{Langer} gives
\begin{equation}\label{fr:Langer-min}
 L_{\min}(V):=
 \lim_{e\to\infty}p^{-e}\mu_{\min}(F_C^{e*}V)
 \ge-\frac{c}{p-1}.
\end{equation}
To make the uniformity explicit, let $g$ be the fixed genus of $C$.
One can choose a nef line bundle $A_C$ of degree at most
$e_g=\max\{4g-2,0\}$ such that $T_C\otimes A_C$ is globally
generated; for $g\le1$, take $A_C=\OO_C$. Langer's bound is
\[
 \mu_{\min}(V)-L_{\min}(V)
 \le\frac{(\rk V-1)\deg A_C}{p-1}.
\]
Since $\rk V=d$ and $\mu_{\min}(V)\ge0$, we may take
$c=(d-1)e_g$, independently of the reduction.

We also need, for every integer $i\ge0$,
\begin{equation}\label{fr:tensor-min}
 \mu_{\min}(V^{\otimes i})\ge iL_{\min}(V).
\end{equation}
Fix $p$ and $i$.  Choose a line bundle $A_e$ of degree
$-\mu_{\min}(F_C^{e*}V)+O(1)$, rounded up with a fixed
genus-dependent additive constant, so that
$F_C^{e*}V\otimes A_e$ is globally generated.  This follows from
Serre duality: after subtracting any point, its minimum slope can
be made greater than $2g(C)-2$, which annihilates $H^1$ and makes
evaluation at that point surjective.  For every vector-bundle
quotient $Q$ of $V^{\otimes i}$, the bundle
$F_C^{e*}Q\otimes A_e^{\otimes i}$ is then globally generated.
Consequently
\[
 p^e\mu(Q)+i\deg A_e\ge0.
\]
Divide by $p^e$ and let $e\to\infty$ to prove
\eqref{fr:tensor-min}.  This limit is taken for each fixed $p,i$;
there is no interchange of Frobenius-iteration and characteristic
limits.

Since $A_{u-j}(V)$ is a quotient of $V^{\otimes(u-j)}$, the perfect
pairing \eqref{fr:socle-pairing} yields
\begin{align*}
 \mu_{\max}(A_j(V))
 &= (p-1)\deg\det V-\mu_{\min}(A_{u-j}(V))\\
 &\le(p-1)K_W\cdot C+\frac{(u-j)c}{p-1}\\
 &\le(p-1)K_W\cdot C+nc.
\end{align*}
\end{proof}

The divisors $B_p-pD=T_0-(p-k_0N_p)G$ range over a fixed finite
list. Combining~\eqref{fr:intersection-bounds},
\eqref{fr:flag-degrees}, and Lemma~\ref{fr:truncated-slope}
therefore gives, uniformly for $0\le j\le u$,
\begin{align}
 \mu_{\max}(\mathcal G_j|_C)
 &\le p(4+2/r)\Lambda H^d+O(1)\notag\\
 &\le M_p:=7p\Lambda H^d+c_0,\label{fr:common-slope-bound}
\end{align}
where $c_0\ge0$ is fixed.  The coefficient $7$ provides harmless
room in this common bound.

\begin{lemma}\label{fr:diagonal-sections}
For all sufficiently large $p$,
\[
 h^0(\Delta_F,\mathcal L_p)
 \le R\bigl(7^nH^d+O(1/p)\bigr)<R/4.
\]
The error constant is independent of the filtration index $j$.
\end{lemma}

\begin{proof}
Write the fixed flag as
\[
 W=W_d\supset W_{d-1}\supset\cdots\supset W_1=C,
 \qquad W_{d-k}\in|h_k|_{W_{d-k+1}}
 \quad(1\le k\le d-1).
\]
At step $k$, fix the preceding nonnegative integers
$b_1,\ldots,b_{k-1}$ and put
\[
 \mathcal F_k=
 \mathcal G_j|_{W_{d-k+1}}
 \Bigl(-\sum_{i<k}b_i h_i\Bigr).
\]
For each $b_k\ge0$, the divisor restriction sequence is
\[
 0\longrightarrow\mathcal F_k(-(b_k+1)h_k)
 \longrightarrow\mathcal F_k(-b_kh_k)
 \longrightarrow
 \mathcal F_k(-b_kh_k)|_{W_{d-k}}
 \longrightarrow0.
\]
Iterating it for $b_k=0,1,\ldots$ and then proceeding to the next
member of the flag bounds the section space by
\begin{equation}\label{fr:flag-lattice-sum}
 h^0(W,\mathcal G_j)\le
 \sum_{b_1,\ldots,b_{d-1}\ge0}
 h^0\left(C,\mathcal G_j|_C\Bigl(-\sum_i b_i h_i|_C\Bigr)\right).
\end{equation}
At each stage the remainder is zero after sufficiently negative
ample twisting.  Its stopping index need not be uniform in $j,p$
or in the preceding twists: enlarging the nonnegative finite sums
to the displayed lattice sum preserves the inequality.

A summand in \eqref{fr:flag-lattice-sum} vanishes when
$\sum_i b_i(h_i\cdot C)>M_p$, because then its maximum slope is
negative.  Each nonzero summand has at most $e_j(M_p+1)$ sections.
Indeed, evaluation at $\lfloor M_p\rfloor+1$ distinct points is injective;
the kernel has negative maximum slope after subtracting those
points.  By \eqref{fr:flag-degrees} it follows that
\begin{align}
 h^0(W,\mathcal G_j)
 &\le e_j(M_p+1)
       \prod_{i=1}^{d-1}
       \left(1+\frac{M_p}{l_i\Lambda H^d}\right)\notag\\
 &=e_jp^d\bigl(7^nH^d+O(1/p)\bigr).
 \label{fr:common-flag-polynomial}
\end{align}
The leading coefficient follows from
$\Lambda=\prod_i l_i$.  More importantly, the whole scalar
polynomial on the first line is the same for every $j$.

Sum over the filtration in \eqref{fr:diagonal-graded}.  The ranks
satisfy $\sum_j e_j=p^d$, so \eqref{fr:common-flag-polynomial} gives
the asserted bound with $R=p^{2d}$.  There is no extra factor for
the number of graded pieces.  Finally
\[
 7^nH^d\le(14\epsilon)^d<1/4
\]
by \eqref{fr:intersection-bounds}, proving the strict inequality
for all large $p$.
\end{proof}

\subsection{The two orders of the determinant}
\label{fr:determinant}

Choose $R$ generically independent columns of
\eqref{fr:two-slot-evaluation} from bases of its source vector spaces.
Their determinant is a nonzero section
\[
 0\ne\sigma\in H^0(W'\times W',\mathcal Q_1\boxtimes\mathcal Q_2)
\]
of an actual external-product line bundle. To specify its factors,
let $m_i$ be the sum of the weights of the $R$ chosen columns in
slot $i$. The target determinant is
$(\det\mathcal E)^s\boxtimes(\det\mathcal E)^s$, and the
chosen source determinant is
$\OO(-m_1S')\boxtimes\OO(-m_2S')$. Thus
\[
 \mathcal Q_i=(\det\mathcal E)^{\otimes s}
                 \otimes\OO_{W'}(m_iS'),\qquad
 \lambda_i=m_i/R.
\]
Every column weight lies between $0$ and $q$, so
$0\le\lambda_i\le q$ independently of the columns chosen.

Under the numerical identification with the base-field twist,
\begin{equation}\label{fr:slot-determinant-class}
 \frac{c_1(\mathcal Q_i)}R
 =\frac{c_1(\mathcal E)}s+\lambda_iS
 =\frac{B_p}p+\frac{p-1}{2p}K_W+\lambda_iS,
 \qquad 0\le\lambda_i\le q.
\end{equation}
The Frobenius first-Chern-class identity is also
\cite[Lemma~4.2]{Sun}; it is an identity of rational divisor classes.  For the second equality, use
\eqref{fr:Frob-pullback-graded}.  Complementary degrees in
\eqref{fr:socle-pairing} have total first Chern class
$e_j(p-1)K_W$.  Summing gives
\[
 \sum_jc_1(A_j(\Omega_W^1))=\frac{s(p-1)}2K_W.
\]
Thus $c_1(F^*\mathcal E)=sB_p+s(p-1)K_W/2$, and Frobenius pullback of
twisted divisor classes multiplies by $p$, as required.

For $0\le\lambda\le q$, the fixed movable inequality gives
\[
 (G+K_W/2+\lambda S)\cdot\gamma\le40\epsilon J\cdot\gamma,
\]
where pullback by $\pi_x$ is understood. This concerns only
intersections of fixed divisors and the fixed complete-intersection
witness, so it remains true on the reduction. The class in
\eqref{fr:slot-determinant-class} differs from its
left-hand divisor at $\lambda=\lambda_i$ by
\[
 \frac1p\bigl(B_p-pD-K_W/2\bigr).
\]
Its numerator belongs to the fixed finite list
$T_0-aD-K_W/2$, $a\in\{0,\ldots,k_0-1\}$. The resulting error after
intersection with $\gamma$ is therefore $O(p^{-1})$, uniformly
over the selected determinant columns.

Let $x'$ be the twisted point and let $0\ne\tau_i\in
H^0(W',\mathcal Q_i)$.  The strict transform of its effective divisor on
the blowup at $x'$ pairs nonnegatively with the twist of $\gamma$.
Since $J\cdot\gamma>0$, it follows that
\begin{equation}\label{fr:slot-order}
 \frac{\ord_{x'}(\tau_i)}R\le40\epsilon+O(1/p).
\end{equation}
The constant is uniform in the columns and in $\tau_i$.  This
argument tests sections on the reduction itself; it does not
require them to lift to characteristic zero.

K\"unneth identifies the space containing $\sigma$ with
$H^0(W',\mathcal Q_1)\otimes H^0(W',\mathcal Q_2)$.  Choose bases adapted to the
orders of vanishing at $x'$ in each factor.  In each degree their
leading terms are linearly independent.  Tensor products of these
leading terms remain independent in each bidegree of the two
disjoint sets of parameters; terms of different bidegrees cannot
cancel.  Every nonzero tensor consequently has order at most the
sum of the two maximum basis orders.  Applying
\eqref{fr:slot-order}, we obtain
\begin{equation}\label{fr:determinant-order-upper}
 \ord_{(x',x')}\sigma\le R\bigl(80\epsilon+O(1/p)\bigr).
\end{equation}

On the other hand, \eqref{fr:diagonal-rank-upper} and
Lemma~\ref{fr:diagonal-sections} bound the matrix rank at $(x',x')$
by a number strictly smaller
than $R/4$.  Trivialize its line bundles near that point.  Invertible
constant row operations make more than $3R/4$ rows vanish at the
point.  Each entry of those rows belongs to the maximal ideal, so
every term of the determinant has order at least $3R/4$.  Therefore
\begin{equation}\label{fr:determinant-order-lower}
 \ord_{(x',x')}\sigma\ge3R/4.
\end{equation}
Equations \eqref{fr:determinant-order-upper} and
\eqref{fr:determinant-order-lower} contradict
$80\epsilon<3/4$ for sufficiently large $p$.


This contradiction proves Proposition~\ref{thm:finite-data}.
The full-rank calculation used $z_*$, while the diagonal estimate
holds at every point and was tested at the separately chosen $x$.
The errors were controlled by one fixed finite list of divisor
classes and one fixed flag polynomial. Thus their constants do
not depend on $p$, the filtration index, or the chosen columns.


\section{Proof of Lemma~\ref{lem:general-nojets}}
\begin{proof}
Fix $s>0$. Choose $n$ sufficiently large that
Lemma~\ref{lem:general-witness} applies, $r_n>s+1$, and
$P_n^d<(2\epsilon)^d$. Since $P_n-r_nD$ is nef,
\[
 (2P_n+sD)\cdot P_n^{d-1}
 \le(2+s/r_n)P_n^d<4P_n^d.
\]
Fix $x,\beta,E,f,H_1,\ldots,H_{d-1}$ supplied by that lemma.
For every $0\le\lambda\le s$, nefness of $D$ gives
\[
 f^*\beta^*(P_n+\lambda D)\cdot H_1\cdots H_{d-1}
 <40\epsilon\,f^*E\cdot H_1\cdots H_{d-1}.
\]
Choose the flag in Lemma~\ref{lem:general-flag} with $P=P_n$.

If $|kR_{n,s}|$ separated $(2d+2)k$-jets at any $z\in U$,
these fixed data would satisfy all hypotheses of
Proposition~\ref{thm:finite-data}, with
\[
 W=X,\qquad G=H=P_n,\qquad S=D,\qquad q=s,\qquad r=r_n.
\]
Here $K_X\sim0$, so the canonical-intersection condition is
automatic. This is impossible. Hence no such $k,z$ exist.
\end{proof}

\medskip
\noindent\textit{Remark.}
This obstruction uses $(\mathsf I_d)$ and $D^d=0$, but does not
use $\kappa(D)=-\infty$. That extra assumption retains the
nonvanishing proof format of the threefold framework. Under
the source inputs, the two jet lemmas actually exclude the
entire full-nef-dimension, non-big case.

\begin{thebibliography}{99}
\bibitem{NS}
\emph{Numerical Semiampleness of Nef Adjoint Divisors},
OpenAI mathematics collection, October 3, 2026,
Sections~2--6, especially Proposition~5.3 and Theorem~6.1.
\href{https://github.com/openai/math/blob/main/preprints/Numerical-Semiampleness-of-Nef-Adjoint-Divisors-October-3-2026/paper.pdf}{Source manuscript}.

\bibitem{BZ}
C.~Birkar and D.-Q.~Zhang,
\emph{Effectivity of Iitaka fibrations and pluricanonical systems of polarized pairs},
Publ. Math. Inst. Hautes \'Etudes Sci. \textbf{123} (2016), 283--331.
\href{https://doi.org/10.1007/s10240-016-0080-x}{DOI}.

\bibitem{BDPP}
S.~Boucksom, J.-P.~Demailly, M.~P\u{a}un, and Th.~Peternell,
\emph{The pseudo-effective cone of a compact K\"ahler manifold and varieties of negative Kodaira dimension},
J.~Algebraic Geom. \textbf{22} (2013), no.~2, 201--248.
\href{https://doi.org/10.1090/S1056-3911-2012-00574-8}{DOI}.

\bibitem{CP}
F.~Campana and M.~P\u{a}un,
\emph{Orbifold generic semi-positivity: an application to families of canonically polarized manifolds},
Ann.~Inst.~Fourier (Grenoble) \textbf{65} (2015), no.~2, 835--861.
\href{https://doi.org/10.5802/aif.2945}{DOI}.

\bibitem{FG}
O.~Fujino and Y.~Gongyo,
\emph{On the moduli $b$-divisors of lc-trivial fibrations},
Ann. Inst. Fourier (Grenoble) \textbf{64} (2014), no.~4, 1721--1735.
\href{https://doi.org/10.5802/aif.2894}{DOI}.

\bibitem{FH}
O.~Fujino and K.~Hashizume,
\emph{Existence of log canonical modifications and its applications},
Eur. J. Math. \textbf{9} (2023), no.~1, Paper No.~13, 32~pp.
\href{https://doi.org/10.1007/s40879-023-00598-0}{DOI}.

\bibitem{KS}
Y.~Kitadai and H.~Sumihiro,
\emph{Canonical filtrations and stability of direct images by Frobenius morphisms},
Tohoku Math. J. (2) \textbf{60} (2008), no.~2, 287--301.
\href{https://doi.org/10.2748/tmj/1215442876}{DOI}.

\bibitem{Langer}
A.~Langer,
\emph{Semistable sheaves in positive characteristic},
Ann. of Math. (2) \textbf{159} (2004), no.~1, 251--276.
\href{https://doi.org/10.4007/annals.2004.159.251}{DOI}.

\bibitem{LazarsfeldPAGI}
R.~Lazarsfeld,
\emph{Positivity in Algebraic Geometry I: Classical Setting: Line Bundles and Linear Series},
Ergebnisse der Mathematik und ihrer Grenzgebiete, vol.~48,
Springer-Verlag, 2004.
\href{https://www.math.stonybrook.edu/~roblaz/Reprints/Laz.PAG.Chapt1.pdf}{Author-hosted Chapter~1}.

\bibitem{MR}
V.~B.~Mehta and A.~Ramanathan,
\emph{Semistable sheaves on projective varieties and their restriction to curves},
Math.~Ann. \textbf{258} (1982), no.~3, 213--224.
\href{https://doi.org/10.1007/BF01450677}{DOI}.

\bibitem{MustataSchwede}
M.~Musta\c{t}\u{a} and K.~Schwede,
\emph{A Frobenius variant of Seshadri constants},
Math. Ann. \textbf{358} (2014), nos.~3--4, 861--878.
\href{https://doi.org/10.1007/s00208-013-0976-4}{DOI}.
Numbering follows arXiv:1203.1081v1.

\bibitem{LA}
OpenAI,
\emph{Log abundance in characteristic zero},
OpenAI Math Release preprint
\href{https://github.com/openai/math/blob/main/preprints/Log-abundance-in-characteristic-zero-September-24-2026/paper.pdf}{OAI:Log-abundance-in-characteristic-zero-September-24-2026},
2026.

\bibitem{Sun}
X.~Sun,
\emph{Direct images of bundles under Frobenius morphism},
Invent. Math. \textbf{173} (2008), no.~2, 427--447.
\href{https://doi.org/10.1007/s00222-008-0125-y}{DOI}.
Numbering follows arXiv:math/0611360v2.
\end{thebibliography}

\end{document}
